\chapter{Additional Problems from MOP 2012}
\begin{quote}
  If our solution has at table of contents, how should we number the pages?
  \\ -- Evan Chen, MOP 2012
\end{quote}

Here's a more recent collection of geometry problems solvable using barycentric coordinates.  These are all problems which appeared at MOP 2012,\footnote{MOP, occasionally called MOSP, stands for Mathematical Olympiad Program; it is the USA training program for the IMO.} along with an IMO question from that year.

These are in very roughly ascending order of difficulty, based both on the amount of computation and the difficulty of setup.
\begin{enumerate}
  \ii (IMO 2012) Given triangle $ABC$ the point $J$ is the center of the excircle opposite the vertex $A$. This excircle is tangent to the side $BC$ at $M$, and to the lines $AB$ and $AC$ at $K$ and $L$, respectively. The lines $LM$ and $BJ$ meet at $F$, and the lines $KM$ and $CJ$ meet at $G.$ Let $S$ be the point of intersection of the lines $AF$ and $BC$, and let $T$ be the point of intersection of the lines $AG$ and $BC.$ Prove that $M$ is the midpoint of $ST.$
  \ii (USA TST 2012) In acute triangle $ABC$, $\angle A < \angle B$ and $\angle A < \angle C$.  Let $P$ be a variable point on side $BC$.  Points $D$ and $E$ lie on sides $AB$ and $AC$, respectively, such that $BP = PD$ and $CP = PE$.  Prove that if $P$ moves along side $BC$, the circumcircle of triangle $ADE$ passes through a fixed point other than $A$.
  \ii (ELMO\footnote{This year's test was called ``Every Little Mistake $\implies$ 0''.} 2012/5) Let $ABC$ be an acute triangle with $AB < AC$, and let $D$ and $E$ be points on side $BC$ such that $BD$ = $CE$ and $D$ lies between $B$ and $E$.  Suppose there exists a point $P$ inside $ABC$ such that $PD \parallel AE$ and $\angle PAB = \angle EAC$.  Prove that $\angle PBA = \angle PCA$.
  \ii (Saudi Arabia TST 2012) Let $ABCD$ be a convex quadrilateral such that $AB = AC = BD$.  The lines $AC$ and $BD$ meet at point $O$, the circles $ABC$ and $ADO$ meet again at point $P$, and the lines $AP$ and $BC$ meet at point $Q$.  Show that $\angle COQ = \angle DOQ$.
  \ii (ISL 2011/G2)  Let $A_1A_2A_3A_4$ be a non-cyclic quadrilateral.  For $1 \le i \le 4$, let $O_i$ and $r_i$ be the circumcenter and the circumradius of triangle $A_{i+1}A_{i+2}A_{i+3}$ (where $A_{i+4} = A_i$).  Prove that
    \[ \frac{1}{O_1A_1^2-r_1^2} + \frac{1}{O_2A_2^2-r_2^2} + \frac{1}{O_3A_3^2-r_3^2} + \frac{1}{O_4A_4^2-r_4^2} = 0 \]
  \ii (USA TSTST 2012) Triangle $ABC$ is inscribed in circle $\Omega$.  The interior angle bisector of angle $A$ intersects side $BC$ and $\Omega$ at $D$ and $L$ (other than $A$), respectively.  Let $M$ be the midpoint of side $BC$.  The circumcircle of triangle $ADM$ intersects sides $AB$ and $AC$ again at $Q$ and $P$ (other than $A$), respectively.  Let $N$ be the midpoint of segment $PQ$, and let $H$ be the foot of the perpendicular from $L$ to line $ND$.  Prove that line $ML$ is tangent to the circumcircle of triangle $HMN$.
  \ii (ISL 2011/G4) Acute triangle $ABC$ is inscribed in circle $\Omega$.  Let $B_0$, $C_0$, $D$ lie on sides $AC$, $AB$, $BC$, respectively, such that $AB_0 = CB_0$, $AC_0 = BC_0$, and $AD \perp BC$.  Let $G$ denote the centroid of triangle $ABC$, and let $\omega$ be a circle through $B_0$ and $C_0$ that is tangent to $\Omega$ at a point $X$ other than $A$.  Prove that $D$, $G$, $X$ are collinear.
  \ii Let $ABC$ be a triangle with circumcenter $O$ and let the angle bisector of $\angle BAC$ intersect $BC$ at $D$.  The point $M$ is such that $MC \perp BC$ and $MA \perp AD$.  Lines $BM$ and $OA$ intersect at the point $P$.  Show that the circle centered at $P$ and passing through $A$ is tangent to segment $BC$.
  \ii (USA TSTST 2012) Let $ABCD$ be a quadrilateral with $AC = BD$.  Diagonals $AC$ and $BD$ meet at $P$.  Let $\omega_1$ and $O_1$ denote the circumcircle and circumcenter of triangle $ABP$.  Let $\omega_2$ and $O_2$ denote the circumcircle and circumcenter of triangle $CDP$.  Segment $BC$ meets $\omega_1$ and $\omega_2$ again at $S$ and $T$ (other than $B$ and $C$), respectively.  Let $M$ and $N$ be the midpoints of minor arcs $SP$ (not including $B$) and $TP$ (not including $C$).  Prove that $MN \parallel O_1O_2$.
\end{enumerate}
